\section{K-Means \& DBSCAN Clustering}

The Python part of this exercise is available in \myhref{https://github.com/LosDelDGIIM/LosDelDGIIM.github.io/blob/main/subjects/Erasmus-DUE/GKI/Exercises/Exercise5.ipynb}{this Jupyter Notebook}.\\

\begin{ejercicio}
    What is the main mathematical difference between K-means and K-medians? Which algorithm is more robust against outliers?

    The main mathematical difference between K-means and K-medians lies in the way they calculate the cluster centers. K-means uses the mean (average) of the data points in a cluster to determine the cluster center, while K-medians uses the median (the middle value when the data points are sorted) to determine the cluster center. The K-medians algorithm is more robust against outliers because the median is less affected by extreme values than the mean. In K-means, outliers can significantly skew the cluster centers, while in K-medians, the median will not be influenced as much by outliers, leading to more stable cluster centers.
\end{ejercicio}

\begin{ejercicio}
    Explain the core idea behind DBSCAN. What do the parameters \texttt{eps} ($\veps$) and $\texttt{min\_samples}$ mean?

    DBSCAN (Density-Based Spatial Clustering of Applications with Noise) is a clustering algorithm that groups together points that are closely packed together, while marking points that lie alone in low-density regions as outliers. The core idea behind DBSCAN is to identify clusters based on the density of data points in the feature space. The two main parameters of DBSCAN are:
    \begin{itemize}
        \item \texttt{eps} ($\veps$): This parameter defines the maximum distance between two points for them to be considered as part of the same cluster. It is essentially the radius of the neighborhood around a point.
        \item $\texttt{min\_samples}$: This parameter specifies the minimum number of points required to form a dense region (i.e., a cluster). A point is considered a core point if it has at least $\texttt{min\_samples}$ points within its $\veps$ neighborhood.
    \end{itemize}
\end{ejercicio}

\begin{ejercicio}
    In the context of DBSCAN, define the terms: core point, border point, and noise (outlier).
    \begin{itemize}
        \item Core Point: A core point is a point that has at least $\texttt{min\_samples}$ points (including itself) within its $\veps$ neighborhood. Core points are the central points of clusters.
        \item Border Point: A border point is a point that is not a core point but is within the $\veps$ neighborhood of a core point. Border points are part of a cluster but do not have enough neighboring points to be considered core points themselves.
        \item Noise (Outlier): A noise point, or outlier, is a point that is neither a core point nor a border point. It does not belong to any cluster and is considered an anomaly in the dataset.
    \end{itemize}
\end{ejercicio}

\begin{ejercicio}
    When does K-Medians (or K-Means) fail at clustering, and when does DBSCAN?

    K-Medians (or K-Means) can fail at clustering when the data has non-spherical clusters, varying cluster sizes, or when there are outliers present in the dataset. These algorithms assume that clusters are spherical and of similar size, which may not always be the case in real-world data. Additionally, K-Means is sensitive to outliers, which can skew the cluster centers and lead to poor clustering results.
\end{ejercicio}